% Resumen analítico (máximo 500 palabras) y su versión en inglés.
\section*{Resumen}

En una red descentralizada no existe un equipo central que vea o dirija todo el
sistema: cada nodo decide con lo que alcanza a oír de sus vecinos. Para estudiar
una red así antes de construirla se usan simuladores, pero las herramientas
disponibles dejan vacíos. Los simuladores de propósito general evalúan una copia
del protocolo escrita para la herramienta; los emuladores ejecutan el programa
real, pero cada nodo consume muchos recursos y no incluyen un modelo listo de la
pérdida de señal entre los pisos de un edificio, y las herramientas de redes
entre pares no representan el medio de radiofrecuencia. Este trabajo implementa
un simulador que atiende esos vacíos. Está escrito en Python y ejecuta en cada
nodo simulado las mismas clases de enrutamiento B.A.T.M.A.N. y de detección de
fallos que usa el nodo de red; lo único que reemplaza es la red, por un modelo
del medio que depende de la distancia y de los pisos. El desarrollo siguió cuatro
etapas: requerimientos, diseño, implementación y pruebas. El simulador envía
paquetes de datos salto a salto con la tabla de rutas de cada nodo, permite elegir
el tipo de paquete (UDP, ICMP o TCP), mueve los Gateway solo con información
local, programa caídas de nodos, repite experimentos con semillas y se maneja
desde una interfaz web local.

Las pruebas son simulaciones de una Red de Expansión de Cobertura con 7 Gateway
y 6 Nodos de usuario en un edificio de tres pisos, con 10 semillas por prueba. La
caída de un Gateway puente mantuvo la malla partida 90 s sin reducir la cobertura
media (1.000): la cobertura, por sí sola, no muestra el daño. Dos fallos dirigidos
al nodo más crítico redujeron la entrega de datos a 0.284 frente a 0.502 de dos
fallos al azar (diferencia pareada de 0.218, IC 95~\% [0.149, 0.287]), con una
cobertura estadísticamente indistinguible. Con Gateway móviles la malla se adaptó
sola: la partición por la caída del puente bajó de 90 s a 6.2 s, a costa de que la
cobertura media descendiera a 0.922. Con el mismo tráfico, TCP entregó 0.924 de
los paquetes, UDP 0.772 e ICMP 0.646, y TCP necesitó cerca de 2.6 veces las
transmisiones de UDP. La inundación entregó más que B.A.T.M.A.N. en campo
abierto, con cerca del doble de transmisiones, pero mucho menos en el edificio
(0.430 contra 0.786). Las tablas de resultados se reprodujeron byte a byte en dos
máquinas. El simulador no se validó contra equipos físicos ni contra el módulo
batman-adv, y su modelo del medio es simple: los resultados describen esta
implementación bajo este medio. Esos frentes quedan como trabajo futuro.

\textbf{Palabras clave}: redes descentralizadas, simulación, B.A.T.M.A.N., redes
malladas, cobertura, adaptabilidad, movilidad descentralizada, reproducibilidad.

% Resumen en inglés
\newpage
\section*{Abstract}

In a decentralized network no central device sees or controls the whole system:
each node decides based on what it hears from its neighbours. Simulators are used
to study such networks before building them, but the available tools leave gaps.
General-purpose simulators evaluate a copy of the protocol written for the tool;
emulators run the real program, but each node is costly and they offer no
ready-made model of signal loss between the floors of a building; and
peer-to-peer tools do not represent the radio-frequency medium. This work
implements a simulator that addresses those gaps. It is written in Python and
runs, in every simulated node, the same B.A.T.M.A.N. routing and
failure-detection classes used by the network node; the only thing it replaces is
the network, with a medium model based on distance and floors. Development
followed four stages: requirements, design, implementation and testing. The
simulator forwards data packets hop by hop using each node's own routing table,
lets the user choose the packet type (UDP, ICMP or TCP), moves Gateways using
local information only, schedules node failures, repeats experiments with seeds
and is operated from a local web interface.

The tests are simulations of a Coverage Expansion Network with 7 Gateways and 6
user nodes in a three-storey building, with 10 seeds per test. Losing a bridge
Gateway kept the mesh partitioned for 90 s without lowering mean coverage (1.000):
coverage alone does not show the damage. Two failures targeting the most critical
node reduced data delivery to 0.284 against 0.502 for two random failures (paired
difference 0.218, 95\% CI [0.149, 0.287]), with statistically indistinguishable
coverage. With mobile Gateways the mesh adapted by itself: the partition caused
by the bridge failure fell from 90 s to 6.2 s, at the cost of mean coverage
falling to 0.922. Under the same traffic, TCP delivered 0.924 of the packets, UDP
0.772 and ICMP 0.646, and TCP needed about 2.6 times the transmissions of UDP.
Flooding delivered more than B.A.T.M.A.N. in open terrain, with about twice the
transmissions, but far less in the building (0.430 against 0.786). Result tables
were reproduced byte for byte on two machines. The simulator was not validated
against physical devices or the batman-adv module, and its medium model is
simple: the results describe this implementation under this medium. Those fronts
remain as future work.

\textbf{Keywords}: decentralized networks, simulation, B.A.T.M.A.N., mesh networks,
coverage, adaptability, decentralized mobility, reproducibility.
